% Generated draft: requires mathematical and source review.
\documentclass[11pt]{article}
\usepackage[margin=25mm]{geometry}
\usepackage{fontspec}
\setmainfont{Latin Modern Roman}
\usepackage{amsmath,amssymb,mathtools,unicode-math}
\setmathfont{Latin Modern Math}
\usepackage{enumitem,needspace,xurl,hyperref}
\setlength{\parindent}{0pt}
\setlength{\parskip}{5pt}
\setlength{\emergencystretch}{3em}
\newcommand{\N}{\mathbb{N}}
\newcommand{\Z}{\mathbb{Z}}
\newcommand{\Q}{\mathbb{Q}}
\newcommand{\R}{\mathbb{R}}
\newcommand{\C}{\mathbb{C}}
\newcommand{\F}{\mathbb{F}}
\newcommand{\A}{\mathbb{A}}
\newcommand{\PP}{\mathbb{P}}
\newcommand{\E}{\mathbb{E}}
\newcommand{\eps}{\varepsilon}
\newcommand{\dd}{\,\mathrm{d}}
\begin{document}
\section*{UnsolvedMath ID: 30003795}
\textbf{Source problem code:} catalogue entry
\subsection{Definitions and mathematical statement}
Let \(K\) be a field finitely generated over its prime subfield. An absolutely almost simple \(K\)-group is a smooth connected affine algebraic \(K\)-group whose base change to an algebraic closure of \(K\) is noncommutative and semisimple and has no proper nontrivial connected normal algebraic subgroup. Fix such a group \(G\).

An inner \(K\)-form of \(G\) is a \(K\)-group \(H\) admitting an isomorphism \(f:H_{K^{\mathrm{sep}}}\to G_{K^{\mathrm{sep}}}\) such that, for every \(\sigma\in\operatorname{Gal}(K^{\mathrm{sep}}/K)\), the automorphism \(f\circ\sigma(f)^{-1}\) of \(G_{K^{\mathrm{sep}}}\) is inner. Forms are counted up to \(K\)-isomorphism of algebraic groups.

For a nontrivial discrete valuation \(v\) of \(K\), write \(\mathcal O_v\) for its valuation ring. Use the standard meaning of good reduction: \(H\) has good reduction at \(v\) if there exists a reductive group scheme \(\mathcal H\) over \(\mathcal O_v\) whose generic fiber is \(K\)-isomorphic to \(H\). Here a reductive group scheme is a smooth affine group scheme of finite presentation with connected reductive geometric fibers.

The supplied clean formulation asks: for every finitely generated field \(K\) of ``good characteristic'' and every absolutely almost simple \(K\)-group \(G\), does there exist a ``natural'' set \(V\) of discrete valuations of \(K\) such that, for every finite subset \(S\subseteq V\),
\[
\#\left\{
[H]_K\;\middle|\;
\begin{array}{l}
H\text{ is an inner }K\text{-form of }G,\\
H\text{ has good reduction at every }v\in V\setminus S
\end{array}
\right\}<\infty?
\]
The bound may depend on \(K,G,V,S\); the statement supplies no uniform numerical bound.

This is a conditional statement schema requiring source clarification. The supplied record does not define ``good characteristic'' or specify which sets of valuations qualify as ``natural.'' These qualifications cannot be dropped or replaced by arbitrary choices of valuations. Moreover, the original formulation refers to an unavailable ``Question 1'' and to ``relevant \(K\)-forms.'' Its wording may require a single valuation set \(V\), depending only on \(K\), to work for every \(G\), whereas the clean formulation above permits \(V\) to depend on \(K\) and \(G\). The original wording also does not itself specify that the relevant forms are precisely the inner forms. These unresolved points prevent identification of a unique fully specified conjecture from the supplied record alone.
\subsection{Short English statement}
The supplied clean question asks whether, over every finitely generated field of suitable characteristic, each absolutely almost simple group admits a natural valuation set for which only finitely many inner forms have good reduction outside any prescribed finite subset. The characteristic restriction, admissible valuation sets, and possible requirement that one valuation set work for all groups remain unspecified.
\subsection{Sources}
\begin{itemize}
\item Local UnsolvedMath registry, supplied complete source record: permanent ID \(30003795\), title \emph{Good-Reduction Finiteness over Finitely Generated Fields}, exact source problem code \texttt{catalogue entry}; collection: \emph{Oberwolfach Reports Open Problems}. No individual proposer is supplied. The proposed year is recorded as \(2018\).
\item Supplied report citation: \emph{Quadratic Forms and Related Structures over Fields (2019)}, Oberwolfach Reports, pp.\ 1235--1291, DOI \texttt{10.4171/owr/2018/21}. \url{https://doi.org/10.4171/owr/2018/21}. The report title/year and DOI are retained exactly as supplied; their dating has not been independently reconciled.
\item Supplied literature source, labeled ``good-reduction survey'': \url{https://arxiv.org/abs/2005.05484}.
\item Supplied literature source, labeled ``recent good-reduction results'': \url{https://www.sciencedirect.com/science/article/abs/pii/S0001870823005807}.
\item Supplied literature source, labeled ``workshop report'': \url{https://ems.press/content/serial-article-files/46742}.
\item Supplied dated literature curation, checked \texttt{2026-08-21}: status ``partially solved.'' Its assessment says that the conjecture is established for inner type \(A\) and several further classes, while current surveys and later work still treat the all-groups, finitely generated, good-characteristic claim as open. This is supplied dated curation derived from the OWR record, not an independently verified literature assessment; its time-sensitive status and cited claims require independent verification.
\item The registry labels the statement \texttt{corrected\_verified} and reports that it was checked against the official Oberwolfach report, with the clean statement preserving its hypotheses, notation, and open target. This is a supplied verification claim, not verification performed for this draft.
\end{itemize}
\end{document}
